\documentclass{article}
\usepackage{amsmath,amssymb,amsthm}
\usepackage{algorithm2e}
\usepackage{booktabs}
\usepackage{graphicx}
\usepackage{hyperref}
\usepackage{xcolor}
\newtheorem{theorem}{Theorem}
\newtheorem{lemma}{Lemma}
\newtheorem{assumption}{Assumption}
\newtheorem{remark}{Remark}
\newcommand{\E}{\mathbb{E}}
\newcommand{\Var}{\operatorname{Var}}
\newcommand{\norm}[1]{\left\lVert#1\right\rVert}
\newcommand{\inner}[2]{\left\langle #1,#2\right\rangle}
\begin{document}

\section*{Algorithm Name}
\textbf{DP‑SGD with Gaussian Noise (DP‑SGD‑G).}  
At each iteration a Gaussian perturbation is added to the (possibly minibatch) gradient to obtain differential privacy. The perturbed gradient is then used in a standard SGD update.

\section*{Mathematical Formulation}
Let $f:\mathbb{R}^{d}\rightarrow\mathbb{R}$ be the objective function.  
Given a stepsize (learning rate) $\eta>0$, noise scale $\sigma>0$ and total number of iterations $T$, the algorithm generates a sequence $\{x_t\}_{t=0}^{T}$ as follows:
\[
\boxed{
\begin{aligned}
g_t &:= \nabla f(x_t)+\xi_t, \qquad \xi_t\sim\mathcal N(0,\sigma^{2}I_d),\\[4pt]
x_{t+1} &:= x_t-\eta\, g_t ,\qquad t=0,1,\dots ,T-1 .
\end{aligned}}
\]
The random vectors $\{\xi_t\}$ are independent of each other and of the data; they provide $(\varepsilon,\delta)$‑DP after appropriate privacy accounting (e.g., moments accountant).

\section*{Pseudocode}
\begin{algorithm}[H]
\SetAlgoLined
\KwIn{Initial point $x_0\in\mathbb R^{d}$, stepsize $\eta$, noise scale $\sigma$, total iterations $T$}
\For{$t=0$ \KwTo $T-1$}{
    Compute (exact or minibatch) gradient $\nabla f(x_t)$\;
    Sample $\xi_t\sim\mathcal N(0,\sigma^{2}I_d)$\;
    Set $g_t\leftarrow \nabla f(x_t)+\xi_t$\;
    Update $x_{t+1}\leftarrow x_t-\eta g_t$\;
}
\KwOut{Final iterate $x_T$ (or averaged iterate $\bar x_T=\frac1T\sum_{t=0}^{T-1}x_t$)}
\caption{DP‑SGD‑G}
\end{algorithm}

\section*{Dry Run Example}
Consider the simple 1‑dimensional quadratic $f(w)=(w-3)^2$.  
Take $w_0=0$, stepsize $\eta=0.1$, noise variance $\sigma^2=0.01$ (i.e., $\sigma=0.1$), and run three iterations.

\begin{tabular}{c c c c c}
\toprule
$t$ & $w_t$ & $\nabla f(w_t)=2(w_t-3)$ & $\xi_t\sim\mathcal N(0,0.01)$ & $w_{t+1}=w_t-\eta(\nabla f(w_t)+\xi_t)$ \\
\midrule
0 & $0$ & $-6$ & (sample) $0.05$ & $0-0.1(-6+0.05)=0.595$ \\
1 & $0.595$ & $-4.81$ & (sample) $-0.02$ & $0.595-0.1(-4.81-0.02)=1.077$ \\
2 & $1.077$ & $-3.846$ & (sample) $0.12$ & $1.077-0.1(-3.846+0.12)=1.452$ \\
\bottomrule
\end{tabular}

Corresponding objective values:
\[
f(w_0)=9,\quad f(w_1)=(0.595-3)^2\approx5.79,\quad
f(w_2)=(1.077-3)^2\approx3.70,\quad
f(w_3)=(1.452-3)^2\approx2.40.
\]
The noisy updates still decrease the objective, illustrating the algorithm’s behaviour on a convex toy problem.

\newpage
\section*{Assumptions}
\begin{assumption}[Smoothness]
The function $f$ is $L$‑smooth, i.e.
\[
\forall x,y\in\mathbb R^{d}:\quad
f(y)\le f(x)+\inner{\nabla f(x)}{y-x}+\frac{L}{2}\norm{y-x}^{2}.
\]
\end{assumption}
\begin{assumption}[Bounded Gradient Norm]
There exists $G>0$ such that $\norm{\nabla f(x)}\le G$ for all $x$ in the domain explored by the algorithm.
\end{assumption}
\begin{assumption}[Gaussian Noise]
The added noise $\xi_t$ satisfies $\xi_t\sim\mathcal N(0,\sigma^{2}I_d)$, independent across iterations and independent of the data.
\end{assumption}
\begin{assumption}[Stepsize]
The stepsize $\eta$ is constant and chosen so that $0<\eta\le \frac{1}{L}$.
\end{assumption}

\section*{Main Convergence Theorem}
\begin{theorem}[Expected Gradient Norm Decay]\label{thm:main}
Under Assumptions~1–4, after $T$ iterations of DP‑SGD‑G the averaged iterate
\[
\bar x_T:=\frac{1}{T}\sum_{t=0}^{T-1}x_t
\]
satisfies
\[
\frac{1}{T}\sum_{t=0}^{T-1}\E\!\bigl[\norm{\nabla f(x_t)}^{2}\bigr]
\;\le\;
\frac{2\bigl(f(x_0)-f^{\star}\bigr)}{\eta T}
\;+\;
L\eta\sigma^{2}d,
\]
where $f^{\star}=\inf_{x}f(x)$ and the expectation is taken over all noise realizations.
Consequently, choosing $\eta = \Theta\!\bigl(1/\sqrt{T}\bigr)$ yields
\[
\frac{1}{T}\sum_{t=0}^{T-1}\E\!\bigl[\norm{\nabla f(x_t)}^{2}\bigr]
= O\!\Bigl(\frac{1}{\sqrt{T}}\Bigr)+O\!\bigl(\sigma^{2}d/\sqrt{T}\bigr).
\]
\end{theorem}

\section*{Theoretical Properties}
\begin{itemize}
\item \textbf{Stability.} The smoothness bound together with the stepsize restriction $0<\eta\le 1/L$ guarantees that the iterates remain in a bounded region, preventing explosion despite the added noise.
\item \textbf{Hyper‑parameter Sensitivity.} The bound exhibits a linear dependence on $\eta$ for the noise term $L\eta\sigma^{2}d$; larger stepsizes accelerate the deterministic decay term but amplify the variance contribution. The optimal $\eta$ balances these two effects, leading to the $\Theta(1/\sqrt{T})$ choice.
\item \textbf{Comparison to Baselines.}
    \begin{enumerate}
        \item \emph{Plain SGD (no noise).} Setting $\sigma=0$ recovers the classic non‑convex SGD bound $\frac{2(f(x_0)-f^{\star})}{\eta T}$.
        \item \emph{DP‑SGD with clipping only.} When clipping is applied but no explicit Gaussian noise, the variance term disappears; however, clipping introduces bias not captured here.
        \item \emph{Adam‑type adaptive methods.} Adaptive stepsizes can improve practical performance, but their theoretical analysis under DP noise remains an open problem; the current bound provides a baseline for any first‑order method with isotropic Gaussian perturbations.
    \end{enumerate}
\end{itemize}

\section*{Discussion}
The theorem shows that, despite the privacy‑inducing Gaussian perturbation, DP‑SGD‑G retains the same $O(1/\sqrt{T})$ convergence rate (up to an additive term proportional to $\sigma^{2}$). The dependence on the ambient dimension $d$ is inevitable when isotropic Gaussian noise is used; dimensionality reduction or per‑coordinate clipping can mitigate this factor in practice.

\newpage
\section*{Preliminaries (Definitions and Lemmas)}
\begin{lemma}[Smoothness Inequality]\label{lem:smooth}
If $f$ is $L$‑smooth, then for any $x$ and any vector $u$,
\[
f(x-u)\le f(x)-\inner{\nabla f(x)}{u}+\frac{L}{2}\norm{u}^{2}.
\]
\end{lemma}
\begin{proof}
Directly from the definition of $L$‑smoothness with $y=x-u$.
\end{proof}

\begin{lemma}[Noise Moment]\label{lem:noise}
For $\xi\sim\mathcal N(0,\sigma^{2}I_d)$,
\[
\E[\xi]=0,\qquad \E[\norm{\xi}^{2}]=\sigma^{2}d,\qquad 
\E[\inner{\nabla f(x)}{\xi}]=0.
\]
\end{lemma}
\begin{proof}
Standard Gaussian moment properties; independence from $x$ yields zero cross‑term.
\end{proof}

\section*{Main Proof (Step‑by‑Step Derivation)}
We prove Theorem~\ref{thm:main}. All expectations are taken w.r.t. the randomness of $\{\xi_t\}_{t=0}^{T-1}$.

\bigskip
\noindent\textbf{Step 1 (Smoothness expansion).}  
Using Lemma~\ref{lem:smooth} with $u=\eta g_t$,
\begin{align}
f(x_{t+1})
&\le f(x_t)-\inner{\nabla f(x_t)}{\eta g_t}
      +\frac{L}{2}\norm{\eta g_t}^{2}\nonumber\\
&= f(x_t)-\eta\inner{\nabla f(x_t)}{g_t}
      +\frac{L\eta^{2}}{2}\norm{g_t}^{2}. \tag{Step 1}
\end{align}

\noindent\textbf{Step 2 (Insert noisy gradient).}  
Recall $g_t=\nabla f(x_t)+\xi_t$, hence
\begin{align}
\inner{\nabla f(x_t)}{g_t}
&=\inner{\nabla f(x_t)}{\nabla f(x_t)}}
  +\inner{\nabla f(x_t)}{\xi_t}
= \norm{\nabla f(x_t)}^{2}+\inner{\nabla f(x_t)}{\xi_t}. \tag{Step 2}
\end{align}
Similarly,
\begin{align}
\norm{g_t}^{2}
&=\norm{\nabla f(x_t)+\xi_t}^{2}
 =\norm{\nabla f(x_t)}^{2}+2\inner{\nabla f(x_t)}{\xi_t}
   +\norm{\xi_t}^{2}. \tag{Step 3}
\end{align}

\noindent\textbf{Step 3 (Plug back into smoothness bound).}  
Substituting \eqref{Step 2} and \eqref{Step 3} into \eqref{Step 1} gives
\begin{align}
f(x_{t+1})
&\le f(x_t)-\eta\Bigl(\norm{\nabla f(x_t)}^{2}
                     +\inner{\nabla f(x_t)}{\xi_t}\Bigr) \nonumber\\
&\quad +\frac{L\eta^{2}}{2}\Bigl(\norm{\nabla f(x_t)}^{2}
                     +2\inner{\nabla f(x_t)}{\xi_t}
                     +\norm{\xi_t}^{2}\Bigr). \tag{Step 4}
\end{align}

\noindent\textbf{Step 4 (Rearrange terms).}  
Collect the deterministic and stochastic parts:
\begin{align}
f(x_{t+1})
&\le f(x_t)
   -\underbrace{\Bigl(\eta-\frac{L\eta^{2}}{2}\Bigr)}_{\alpha}
      \norm{\nabla f(x_t)}^{2}
   -\underbrace{\Bigl(\eta-\!L\eta^{2}\Bigr)}_{\beta}
      \inner{\nabla f(x_t)}{\xi_t}
   +\frac{L\eta^{2}}{2}\norm{\xi_t}^{2}. \tag{Step 5}
\end{align}
Because $0<\eta\le 1/L$, we have $\alpha\ge \frac{\eta}{2}$ and $\beta\ge 0$.
Thus,
\begin{align}
f(x_{t+1})
&\le f(x_t)-\frac{\eta}{2}\norm{\nabla f(x_t)}^{2}
   -\beta\inner{\nabla f(x_t)}{\xi_t}
   +\frac{L\eta^{2}}{2}\norm{\xi_t}^{2}. \tag{Step 6}
\end{align}

\noindent\textbf{Step 5 (Take expectation).}  
Condition on $x_t$ and use Lemma~\ref{lem:noise}:
\begin{align}
\E\bigl[f(x_{t+1})\mid x_t\bigr]
&\le f(x_t)-\frac{\eta}{2}\norm{\nabla f(x_t)}^{2}
   -\beta\,\E\bigl[\inner{\nabla f(x_t)}{\xi_t}\mid x_t\bigr]
   +\frac{L\eta^{2}}{2}\E\bigl[\norm{\xi_t}^{2}\mid x_t\bigr] \nonumber\\
&= f(x_t)-\frac{\eta}{2}\norm{\nabla f(x_t)}^{2}
   +\frac{L\eta^{2}}{2}\sigma^{2}d. \tag{Step 7}
\end{align}
(Here the cross term vanishes because $\E[\xi_t]=0$.)

\noindent\textbf{Step 6 (Unconditional expectation).}  
Taking total expectation and telescoping over $t=0,\dots,T-1$,
\begin{align}
\E[f(x_T)]
&\le f(x_0)-\frac{\eta}{2}\sum_{t=0}^{T-1}\E\!\bigl[\norm{\nabla f(x_t)}^{2}\bigr]
   +\frac{L\eta^{2}}{2}\sigma^{2}d\,T. \tag{Step 8}
\end{align}
Since $f(x_T)\ge f^{\star}$,
\begin{align}
\frac{\eta}{2}\sum_{t=0}^{T-1}\E\!\bigl[\norm{\nabla f(x_t)}^{2}\bigr]
&\le f(x_0)-f^{\star}
   +\frac{L\eta^{2}}{2}\sigma^{2}d\,T. \tag{Step 9}
\end{align}

\noindent\textbf{Step 7 (Divide by $T$).}
\[
\frac{1}{T}\sum_{t=0}^{T-1}\E\!\bigl[\norm{\nabla f(x_t)}^{2}\bigr]
\le \frac{2\bigl(f(x_0)-f^{\star}\bigr)}{\eta T}
   +L\eta\sigma^{2}d. \tag{Step 10}
\]
This is precisely the bound stated in Theorem~\ref{thm:main}.

\noindent\textbf{Step 8 (Optimizing $\eta$).}  
Setting $\eta = \sqrt{\frac{2\bigl(f(x_0)-f^{\star}\bigr)}{L\sigma^{2}d\,T}}$ (which respects $\eta\le 1/L$ for sufficiently large $T$) yields
\[
\frac{1}{T}\sum_{t=0}^{T-1}\E\!\bigl[\norm{\nabla f(x_t)}^{2}\bigr]
= O\!\Bigl(\frac{1}{\sqrt{T}}\Bigr)
   +O\!\bigl(\sigma^{2}d/\sqrt{T}\bigr),
\]
establishing the claimed $O(1/\sqrt{T})$ rate. \qed

\section*{Rate Derivation (With Constants)}
From \eqref{Step 10} we can write explicitly
\[
\boxed{
\frac{1}{T}\sum_{t=0}^{T-1}\E\!\bigl[\norm{\nabla f(x_t)}^{2}\bigr]
\le
\underbrace{\frac{2\Delta_f}{\eta T}}_{\text{deterministic decay}}
\;+\;
\underbrace{L\eta\sigma^{2}d}_{\text{privacy noise penalty}},
}
\]
where $\Delta_f:=f(x_0)-f^{\star}$.
Choosing $\eta = \min\!\bigl\{1/L,\sqrt{2\Delta_f/(L\sigma^{2}d\,T)}\bigr\}$ gives the optimal trade‑off.

\section*{Special Cases}
\begin{enumerate}
\item \textbf{No noise ($\sigma=0$).} The bound reduces to the classic non‑convex SGD result
\[
\frac{1}{T}\sum_{t=0}^{T-1}\E\!\bigl[\norm{\nabla f(x_t)}^{2}\bigr]
\le \frac{2\Delta_f}{\eta T},
\]
which can be driven to $O(1/T)$ by taking a diminishing stepsize $\eta_t = \Theta(1/t)$.
\item \textbf{Large batch (variance negligible).} If the stochastic gradient already has variance $\tau^{2}$, the total variance becomes $\tau^{2}+\sigma^{2}d$, and the same proof carries through with $\sigma^{2}d$ replaced by $\tau^{2}+\sigma^{2}d$.
\item \textbf{Convex $L$‑smooth case.} Under convexity one can replace the gradient‑norm bound by a function‑value bound, obtaining
\[
\E\bigl[f(\bar x_T)-f^{\star}\bigr]
\le \frac{\Delta_f}{\eta T}+ \frac{L\eta\sigma^{2}d}{2},
\]
which mirrors the standard DP‑SGD convex analysis.
\end{enumerate}

\end{document}